\makeabstract
    {
    Optimizing Activity Assays for Alpha-Lytic Protease: A Foundation for Studying Protein Folding in Lysobacter enzymogenes
    }
    {
    Emma Torres\textsuperscript{1}, Ellen Wolstenholme\textsuperscript{2}, Ryan Tak\textsuperscript{2}, Sheila Jaswal\textsuperscript{1,2}
    }
    {
    \textsuperscript{1} Department of Biology, Amherst College, Amherst, MA\\
\textsuperscript{2} Biochemistry \& Biophysics Program, Amherst College, Amherst, MA
    }
    {
    Purpose: Alpha-lytic protease (\ensuremath{\alpha}LP) and beta-lytic metalloendopeptidase (\ensuremath{\beta}LM) are two enzymes produced by the soil bacteria Lysobacter enzymogenes. Its unique folding process makes both \ensuremath{\alpha}LP and \ensuremath{\beta}LM interesting for industrial and agricultural use. The goal of our study is to develop a method for measuring \ensuremath{\alpha}LP activity.
    }
    {
    Methods: \ensuremath{\alpha}LP activity was measured based on the breakdown of a synthetic substrate (Suc-Ala-Ala-Pro-Phe-pNA). Two measurement techniques were used: a UV/Vis spectrophotometer and a 96-well plate reader, they were compared for their sensitivity, sample volume requirements, and efficiency using serial dilutions of an archived \ensuremath{\alpha}LP stock (2015). In addition to stock, \ensuremath{\alpha}LP L. enzymogenes was cultured. Its growth was studied in two different types of nutrient media (Bachovchin's Media and Tryptone Soya Broth) to evaluate bacterial growth and natural \ensuremath{\alpha}LP secretion, monitored by optical density and activity assays.
    }
    {
    Results: The plate reader produced clean and linear results. It required less volume of sample while allowing for multiple dilutions to be tested simultaneously, making it the more practical method going forward. A 1:100 dilution of the \ensuremath{\alpha}LP stock yielded the most favorable and consistent readings. As for the cultivated  L. enzymogenes bacterial growth differed between the two media, with lower optical density observed in Tryptone Soya Broth. 
    }
    {
    Conclusion: This study establishes a more efficient, plate reader-based protocol for measuring \ensuremath{\alpha}LP activity and identifies key next steps for improving bacterial culture conditions to support future studies. These findings lay the groundwork for extending this method to \ensuremath{\beta}LM and for exploring alternative expression systems, such as Bacillus subtilis, ultimately enabling comparative research into how these two related proteases fold and stabilize.
    }

    \newpage
    